\chapter{Eulerian and Hamiltonian Paths}
\chapterepigraph{Two famous path problems look similar but could not be
more different in difficulty. An Eulerian path uses every \emph{edge}
exactly once and exists under a simple degree condition, findable in
linear time. A Hamiltonian path visits every \emph{node} exactly once, is
NP-hard in general, and here is tamed only by small $n$ (bitmask DP) or a
heuristic on a grid.}

\section{Every Edge Once vs Every Node Once}
\label{sec:cses-graph-euler-intro}

The distinction is the whole chapter: edges (Eulerian, easy) versus nodes
(Hamiltonian, hard).

\begin{patternbox}[Which Path Are You Asked For?]
\begin{itemize}
  \item \textbf{Use every edge exactly once} -- \textbf{Eulerian
        circuit/path} via Hierholzer's algorithm, after checking the
        degree conditions: Mail Delivery (undirected circuit),
        Teleporters Path (directed path), De Bruijn Sequence (Eulerian
        path in a De Bruijn graph).
  \item \textbf{Visit every node exactly once (count or find)} --
        \textbf{Hamiltonian}: bitmask DP over visited sets when $n\le20$
        (Hamiltonian Flights), or a Warnsdorff heuristic on a board
        (Knight's Tour).
  \end{itemize}
\end{patternbox}

\begin{ideabox}[Eulerian Existence Is a Degree Check]
An undirected connected graph has an Eulerian circuit iff every vertex has
even degree; a directed one iff every vertex's in-degree equals its
out-degree (for a path, exactly one $+1/-1$ imbalance pair). Hierholzer's
algorithm then stitches edge-disjoint cycles into a single trail in linear
time.
\end{ideabox}

\newpage

\section{Mail Delivery}
\label{sec:cses-mail-delivery}

\problemstatement{Given an undirected graph of $n$ nodes and $m$ edges,
find a route that starts and ends at node $1$ and traverses \emph{every
edge exactly once} (an Eulerian circuit), or output \texttt{IMPOSSIBLE}.}

\begin{patternbox}
An \textbf{Eulerian circuit} exists iff the graph is connected (over
edge-bearing nodes) and \emph{every vertex has even degree}. When it
exists, \textbf{Hierholzer's algorithm} stitches it together in linear
time.
\end{patternbox}

\subsection*{Degree check, then Hierholzer}

First verify every vertex has even degree and all edges lie in one
connected piece reachable from node $1$; otherwise no circuit exists.
Hierholzer's then walks edges greedily from node $1$, marking each used,
until it returns to a node with no unused edges -- backtracking (popping to
an output list) whenever stuck, which merges sub-cycles into one trail.
Reversing the popped order yields the Eulerian circuit. A per-node pointer
skips already-used edges so total work stays linear.

\begin{ideabox}[Even Degrees Guarantee You Can Always Leave]
Entering a vertex uses one incident edge; an even degree guarantees an
unused edge to leave by, so the walk only gets stuck back at the start.
Hierholzer exploits this to splice every local cycle into a single circuit.
\end{ideabox}

\subsection*{Algorithm}

\begin{enumerate}
  \item Check all degrees even and edges connected; else
        \texttt{IMPOSSIBLE}.
  \item Hierholzer from node $1$ using a stack and per-node edge pointer,
        marking edges used.
  \item Reverse the popped nodes; if it uses all $m$ edges, print it.
\end{enumerate}

\subsection*{Dry run}

\dryrun{Square $1\!-\!2$, $2\!-\!3$, $3\!-\!4$, $4\!-\!1$.}

\begin{itemize}
  \item All degrees $2$ (even), connected -- a circuit exists.
  \item Hierholzer yields $1\to2\to3\to4\to1$.
\end{itemize}

\subsection*{C++ implementation}

\begin{lstlisting}
#include <bits/stdc++.h>
using namespace std;

int main() {
    int n, m;
    cin >> n >> m;
    vector<vector<pair<int,int>>> adj(n + 1);    // {neighbor, edgeId}
    vector<int> deg(n + 1, 0);
    for (int i = 0; i < m; i++) {
        int a, b; cin >> a >> b;
        adj[a].push_back({b, i});
        adj[b].push_back({a, i});
        deg[a]++; deg[b]++;
    }
    for (int i = 1; i <= n; i++)
        if (deg[i] % 2 != 0) { cout << "IMPOSSIBLE\n"; return 0; }

    vector<int> ptr(n + 1, 0);
    vector<bool> used(m, false);
    vector<int> circuit;
    stack<int> st;
    st.push(1);
    while (!st.empty()) {
        int u = st.top();
        while (ptr[u] < (int)adj[u].size() && used[adj[u][ptr[u]].second]) ptr[u]++;
        if (ptr[u] == (int)adj[u].size()) { circuit.push_back(u); st.pop(); }
        else {
            auto [v, id] = adj[u][ptr[u]];
            used[id] = true;
            st.push(v);
        }
    }
    reverse(circuit.begin(), circuit.end());
    if ((int)circuit.size() != m + 1) { cout << "IMPOSSIBLE\n"; return 0; }
    for (int v : circuit) cout << v << " ";
    cout << "\n";
    return 0;
}
\end{lstlisting}

\begin{complexitybox}
\textbf{Time Complexity.} $O(n+m)$ -- each edge is used once; pointers avoid
rescanning. \textbf{Space Complexity.} $O(n+m)$.
\end{complexitybox}

\begin{takeawaybox}
Eulerian circuits are a degree check plus Hierholzer's stitching of cycles.
The per-node pointer and edge-used markers keep it linear even with
multi-edges. The circuit-size check at the end catches disconnection.
\end{takeawaybox}

\section{De Bruijn Sequence}
\label{sec:cses-de-bruijn-sequence}

\problemstatement{For a given $n$, output a binary string that contains
\emph{every} binary string of length $n$ as a substring, and is as short as
possible (length $2^n+n-1$, treating the string as almost cyclic).}

\begin{patternbox}
A De Bruijn sequence is an \textbf{Eulerian path in the De Bruijn graph}:
nodes are the $2^{n-1}$ binary strings of length $n-1$, and each length-$n$
string is an edge from its $(n-1)$-prefix to its $(n-1)$-suffix. Every node
has equal in- and out-degree, so an Eulerian circuit exists.
\end{patternbox}

\subsection*{Edges are the length-n strings}

Build a graph whose nodes are all $(n-1)$-bit strings. For each $n$-bit
string, add an edge from its first $n-1$ bits to its last $n-1$ bits,
labelled by the last bit. Because appending $0$ or $1$ to any node gives a
valid edge, in-degree equals out-degree everywhere -- the Eulerian
condition. Run Hierholzer to traverse every edge once; concatenating the
starting node with each edge's appended bit produces a string in which every
length-$n$ substring appears exactly once.

\begin{ideabox}[Traverse Every Length-n String as an Edge]
Walking an Eulerian circuit uses each length-$n$ string (edge) exactly once,
and consecutive edges overlap in $n-1$ characters, so the emitted bits form
a compact sequence containing all $2^n$ patterns.
\end{ideabox}

\subsection*{Algorithm}

\begin{enumerate}
  \item Nodes $=$ integers $0\ldots2^{n-1}-1$; for each node $v$ and bit
        $b$, edge to $((v\ll1)\,|\,b)\bmod 2^{n-1}$ labelled $b$.
  \item Hierholzer from node $0$, recording edge labels.
  \item Output the start node's $n-1$ bits followed by the recorded labels.
\end{enumerate}

\subsection*{Dry run}

\dryrun{$n=2$: nodes $\{0,1\}$ (1-bit), edges for $00,01,10,11$.}

\begin{itemize}
  \item Eulerian circuit visits all four length-$2$ strings.
  \item A valid output is \texttt{00110}: substrings $00,01,11,10$ all
        appear.
\end{itemize}

\subsection*{C++ implementation}

\begin{lstlisting}
#include <bits/stdc++.h>
using namespace std;

int main() {
    int n;
    cin >> n;
    if (n == 1) { cout << "01\n"; return 0; }

    int nodes = 1 << (n - 1);
    int mask = nodes - 1;
    vector<int> ptr(nodes, 0);                    // next bit to try per node (0 then 1)

    // Hierholzer recording edge labels; stack holds {node, label used to enter}
    vector<int> labels;
    stack<pair<int,int>> st;
    st.push({0, -1});                              // start at node 0, no entering label
    while (!st.empty()) {
        auto [v, lab] = st.top();
        if (ptr[v] < 2) {                          // an unused out-edge (bit 0 then 1)
            int b = ptr[v]++;
            int to = ((v << 1) | b) & mask;
            st.push({to, b});
        } else {
            st.pop();
            if (lab != -1) labels.push_back(lab);
        }
    }
    reverse(labels.begin(), labels.end());

    string res(n - 1, '0');                        // start node = (n-1) zero bits
    for (int b : labels) res += char('0' + b);
    cout << res << "\n";
    return 0;
}
\end{lstlisting}

\begin{complexitybox}
\textbf{Time Complexity.} $O(2^n)$ -- one traversal of the $2^n$ edges.
\textbf{Space Complexity.} $O(2^n)$ for the recursion/circuit.
\end{complexitybox}

\begin{takeawaybox}
A De Bruijn sequence is an Eulerian circuit in the De Bruijn graph, where
length-$n$ strings are edges over $(n-1)$-length nodes. The construction is
the flagship application of Eulerian paths beyond literal edge-traversal
puzzles.
\end{takeawaybox}

\section{Teleporters Path}
\label{sec:cses-teleporters-path}

\problemstatement{Given a \emph{directed} graph of $n$ nodes and $m$ edges,
find a path from node $1$ to node $n$ that uses \emph{every edge exactly
once} (a directed Eulerian path), or output \texttt{IMPOSSIBLE}.}

\begin{patternbox}
A directed \textbf{Eulerian path} from $1$ to $n$ exists iff node $1$ has
one more outgoing than incoming edge, node $n$ one more incoming than
outgoing, every other node has equal in- and out-degree, and all edges are
connected. \textbf{Hierholzer's algorithm} then constructs it.
\end{patternbox}

\subsection*{Degree imbalance pins the endpoints}

Check the directed degree conditions: $\textit{out}[1]-\textit{in}[1]=1$,
$\textit{in}[n]-\textit{out}[n]=1$, and $\textit{in}=\textit{out}$ for all
other nodes; also confirm every edge is reachable so the trail is a single
path. Then run Hierholzer from node $1$: follow unused out-edges via a
per-node pointer, pushing onto a stack; when a node has no unused out-edges,
pop it to the path. Reversing the popped nodes gives the Eulerian path,
which necessarily ends at $n$.

\begin{ideabox}[The +1/-1 Imbalance Fixes Start and End]
Equal in/out degree everywhere would force a circuit; a single node with an
extra out-edge must be the start and a single node with an extra in-edge the
end. Hierholzer respects this automatically, terminating the trail at $n$.
\end{ideabox}

\subsection*{Algorithm}

\begin{enumerate}
  \item Compute in/out degrees; verify the $1\to n$ Eulerian-path
        conditions, else \texttt{IMPOSSIBLE}.
  \item Hierholzer from node $1$ over directed out-edges with a per-node
        pointer.
  \item Reverse the popped nodes; if all $m$ edges are used and it ends at
        $n$, print it.
\end{enumerate}

\subsection*{Dry run}

\dryrun{$1\!\to\!2$, $2\!\to\!3$, $1\!\to\!3$? (degrees must fit).}

\begin{itemize}
  \item With $1\!\to\!2$, $2\!\to\!3$: $\textit{out}[1]{-}\textit{in}[1]{=}
        1$, $\textit{in}[3]{-}\textit{out}[3]{=}1$, node $2$ balanced.
  \item Path $1\to2\to3$ uses both edges.
\end{itemize}

\subsection*{C++ implementation}

\begin{lstlisting}
#include <bits/stdc++.h>
using namespace std;

int main() {
    int n, m;
    cin >> n >> m;
    vector<vector<int>> adj(n + 1);
    vector<int> indeg(n + 1, 0), outdeg(n + 1, 0);
    for (int i = 0; i < m; i++) {
        int a, b; cin >> a >> b;
        adj[a].push_back(b);
        outdeg[a]++; indeg[b]++;
    }

    // Eulerian path 1 -> n conditions
    bool ok = (outdeg[1] - indeg[1] == 1) && (indeg[n] - outdeg[n] == 1);
    for (int i = 1; i <= n && ok; i++)
        if (i != 1 && i != n && indeg[i] != outdeg[i]) ok = false;
    if (!ok) { cout << "IMPOSSIBLE\n"; return 0; }

    vector<int> ptr(n + 1, 0), path;
    stack<int> st;
    st.push(1);
    while (!st.empty()) {
        int u = st.top();
        if (ptr[u] < (int)adj[u].size()) {
            int v = adj[u][ptr[u]++];
            st.push(v);
        } else { path.push_back(u); st.pop(); }
    }
    reverse(path.begin(), path.end());

    if ((int)path.size() != m + 1 || path.front() != 1 || path.back() != n) {
        cout << "IMPOSSIBLE\n"; return 0;
    }
    for (int v : path) cout << v << " ";
    cout << "\n";
    return 0;
}
\end{lstlisting}

\begin{complexitybox}
\textbf{Time Complexity.} $O(n+m)$. \textbf{Space Complexity.} $O(n+m)$.
\end{complexitybox}

\begin{takeawaybox}
Directed Eulerian paths need the in/out degree imbalance to name the
endpoints, then Hierholzer. The final check -- all edges used and the trail
runs $1\to n$ -- guards against disconnected edge sets that pass the local
degree test.
\end{takeawaybox}

\section{Hamiltonian Flights}
\label{sec:cses-hamiltonian-flights}

\problemstatement{Given a directed graph of $n$ cities ($n\le20$) and $m$
flights, count the number of routes from city $1$ to city $n$ that visit
\emph{every} city exactly once (Hamiltonian paths). Output the count modulo
$10^9+7$.}

\begin{patternbox}
Hamiltonian path counting is NP-hard in general, but $n\le20$ invites
\textbf{bitmask DP}: the state is (set of visited cities, current city), and
$2^n\cdot n$ states are affordable.
\end{patternbox}

\subsection*{DP over visited subsets}

Let $\textit{dp}[\textit{mask}][v]$ be the number of routes that start at
city $1$, have visited exactly the set $\textit{mask}$ (which must include
$1$ and $v$), and currently sit at $v$. Base case:
$\textit{dp}[\{1\}][1]=1$. Transition: from state $(\textit{mask},u)$, for
each flight $u\to v$ with $v\notin\textit{mask}$, add
$\textit{dp}[\textit{mask}][u]$ to
$\textit{dp}[\textit{mask}\cup\{v\}][v]$. The answer is
$\textit{dp}[\text{full}][n]$ -- all cities visited, ending at $n$.

\begin{ideabox}[The Visited Set Is the State]
A Hamiltonian path cannot revisit a city, so the exact set already visited
(plus where you are) fully determines valid continuations. Encoding that set
as a bitmask makes the DP finite: $2^n$ subsets times $n$ endpoints.
\end{ideabox}

\subsection*{Algorithm}

\begin{enumerate}
  \item $\textit{dp}[1\ll0][1]=1$ (city $1$ is bit $0$).
  \item For each mask (increasing) and each city $u$ in it with
        $\textit{dp}>0$: extend along flights $u\to v$ with $v$ unvisited.
  \item Output $\textit{dp}[(1\ll n)-1][n-1]$.
\end{enumerate}

\subsection*{Dry run}

\dryrun{$n=3$; flights $1\!\to\!2$, $2\!\to\!3$, $1\!\to\!3$.}

\begin{itemize}
  \item Only $1\to2\to3$ visits all three ending at $3$; $1\to3$ skips
        city $2$.
  \item Count $=1$.
\end{itemize}

\subsection*{C++ implementation}

\begin{lstlisting}
#include <bits/stdc++.h>
using namespace std;
const long long MOD = 1e9 + 7;

int main() {
    int n, m;
    cin >> n >> m;
    vector<vector<int>> adj(n);                    // 0-indexed cities
    for (int i = 0; i < m; i++) {
        int a, b; cin >> a >> b;
        adj[a - 1].push_back(b - 1);
    }

    int full = 1 << n;
    vector<vector<long long>> dp(full, vector<long long>(n, 0));
    dp[1][0] = 1;                                  // visited {city 0}, at city 0
    for (int mask = 1; mask < full; mask++)
        for (int u = 0; u < n; u++) {
            if (!(mask & (1 << u)) || dp[mask][u] == 0) continue;
            for (int v : adj[u])
                if (!(mask & (1 << v)))
                    dp[mask | (1 << v)][v] =
                        (dp[mask | (1 << v)][v] + dp[mask][u]) % MOD;
        }
    cout << dp[full - 1][n - 1] << "\n";
    return 0;
}
\end{lstlisting}

\begin{complexitybox}
\textbf{Time Complexity.} $O(2^n\cdot n^2)$ in the worst case ($n$ endpoints,
$n$ neighbours). \textbf{Space Complexity.} $O(2^n\cdot n)$ for the DP table.
\end{complexitybox}

\begin{takeawaybox}
Hamiltonian counting is bitmask DP with state (visited set, current node);
small $n$ is what makes the exponential state space tractable. This
subset-DP template also solves the travelling salesman problem and
assignment-style problems.
\end{takeawaybox}

\section{Knight's Tour}
\label{sec:cses-knights-tour}

\problemstatement{On an $8\times8$ chessboard, starting from a given square,
find a knight's tour: a sequence of knight moves visiting every square
exactly once. Output the board with each square labelled by the step at
which the knight visits it.}

\begin{patternbox}
A knight's tour is a Hamiltonian path on the knight-move graph. Rather than
exponential search, \textbf{Warnsdorff's heuristic} makes it near-instant:
always move to the square with the \emph{fewest} onward unvisited moves.
\end{patternbox}

\subsection*{Warnsdorff's rule}

From the current square, consider all legal knight moves to unvisited
squares. Move to the one whose \emph{onward degree} -- the number of
unvisited squares \emph{it} can reach -- is smallest. The intuition:
squares with few exits are the easiest to get stranded in, so visit them
early while they are still reachable. Ties can be broken arbitrarily (or by
a fixed order). On a standard board this greedy choice completes a full tour
without backtracking in practice.

\begin{ideabox}[Escape the Tightest Corners First]
Choosing the least-accessible next square keeps the most flexible squares
for later, avoiding dead ends. This ``most-constrained first'' heuristic
turns an exponential Hamiltonian search into an $O(64)$ greedy walk.
\end{ideabox}

\subsection*{Algorithm}

\begin{enumerate}
  \item Mark the start square as step $1$.
  \item Repeat $63$ times: among unvisited knight-move targets, pick the
        one with the fewest unvisited onward moves; go there, label the
        next step.
  \item Output the $8\times8$ grid of step labels.
\end{enumerate}

\subsection*{Dry run}

\dryrun{Start at a corner, e.g.\ $(0,0)$.}

\begin{itemize}
  \item The corner has only two moves; Warnsdorff steers toward low-degree
        squares first, threading all $64$ cells.
  \item The result is a valid tour numbered $1\ldots64$.
\end{itemize}

\subsection*{C++ implementation}

\begin{lstlisting}
#include <bits/stdc++.h>
using namespace std;

int dr[] = {-2,-2,-1,-1,1,1,2,2};
int dc[] = {-1,1,-2,2,-2,2,-1,1};

int main() {
    int sr, sc;
    cin >> sc >> sr;                              // column, row (1-indexed)
    sr--; sc--;

    vector<vector<int>> board(8, vector<int>(8, 0));
    auto inb = [](int r, int c) { return r >= 0 && r < 8 && c >= 0 && c < 8; };
    auto onward = [&](int r, int c) {             // unvisited onward moves
        int cnt = 0;
        for (int k = 0; k < 8; k++) {
            int nr = r + dr[k], nc = c + dc[k];
            if (inb(nr, nc) && board[nr][nc] == 0) cnt++;
        }
        return cnt;
    };

    int r = sr, c = sc;
    board[r][c] = 1;
    for (int step = 2; step <= 64; step++) {
        int best = -1, bestR = -1, bestC = -1;
        for (int k = 0; k < 8; k++) {
            int nr = r + dr[k], nc = c + dc[k];
            if (!inb(nr, nc) || board[nr][nc]) continue;
            int deg = onward(nr, nc);
            if (best == -1 || deg < best) { best = deg; bestR = nr; bestC = nc; }
        }
        r = bestR; c = bestC;
        board[r][c] = step;
    }

    for (int i = 0; i < 8; i++) {
        for (int j = 0; j < 8; j++) cout << board[i][j] << " ";
        cout << "\n";
    }
    return 0;
}
\end{lstlisting}

\begin{complexitybox}
\textbf{Time Complexity.} $O(64)$ moves, each evaluating a constant number
of candidates -- effectively constant for the fixed board. \textbf{Space
Complexity.} $O(1)$ (a fixed $8\times8$ board).
\end{complexitybox}

\begin{takeawaybox}
Warnsdorff's ``fewest onward moves'' heuristic finds a knight's tour without
backtracking by always escaping the most constrained square first. It is the
practical alternative when a full Hamiltonian search would be too slow.
\end{takeawaybox}

